\scanpage{23}

% L1–6
d) \uline{Passage à la limite} Cela correspond encore : des 2-Limites \inserted{inductives} \unsure{dans|\deleted{dans}} $\Delta$, il faut montrer que les 2-limites transfinies existent et sont des topos. La démonstration donnée dans c, marche en fait pour tt système inductif \deleted{de topos} dans $\Delta$, formé de topos [du système projectif dans (Top)] : $(R_\alpha)$ : On se débrouille pour choisir des
% L7–11
\deleted{systèmes de} \inserted{petites cat.} généra\deleted{trices} \inserted{$C_\alpha\subset R_\alpha$} \unsure{qui|\deleted{qui}} sont stables par les morphismes de transition \deleted{\unsure{ce qui est f}} -- p.ex. à l'aide des filtrations cardinales, mais c'est un marteau pour une mouche -- et on prend \deleted{la catégorie-\unsure{limite in|2-} de} $C=\underset{\alpha\,(\mathrm{Cat})}{\underline{\mathrm{Lim}}}\, C_\alpha$, on \deleted{prend} trouve un sous-topos de $\widehat C$.

% L12–17
En l'occurrence, \deleted{on voit que si on} \inserted{dans le cas filtrant} on voit que si \deleted{un} $R$ est le topos limite, \deleted{et si} les images dans $R$ des $R_\alpha$ forment un syst. de générateurs ; \deleted{\unsure{hyp si}} \inserted{\unsure{donc|alors|d'où}} on a une \deleted{syst. de générateurs} catégorie $\mathbf J$ s'envoyant dans le syst. inductif des $R_\alpha$, $J-R_\alpha$ engendre $R_\alpha$ pour tt $\alpha$, alors $J-R$ engendre $R$.

% L18, horizontal rule
\par\noindent\rule{\linewidth}{.3pt}\par

% L18–24, lower diagram
\diagram{\begin{tikzcd}[ampersand replacement=\&,column sep=small,row sep=small]
 \rho(I)=\widehat J_0\arrow[d,equal] \& \rho'(J_1)\arrow[d,equal]\& \rho^{\prime}(J_2)\arrow[d,equal]\& {}\\
 R_0\arrow[r] \& R_1\arrow[r,"\mathrm{loc}"]\& R_2\arrow[r,"\mathrm{loc}"]\& R_3=R\\
 J_0\arrow[u]\arrow[r,"\substack{\text{bij sur}\\\text{objets}}"]\&
 J_1\arrow[u]\arrow[r,"\mathrm{loc.}","{(\text{bij sur obj})}"']\&
 J_2\arrow[u]\arrow[ur,leftrightarrow]\& {}\\
 I\arrow[u]\& {} \& {} \& {}
\end{tikzcd}}

% L18, annotation above diagram
$\rho'$ la $\Delta_\pi$-\unsure{b|\Delta}-enveloppe, où $\Delta_\pi$ cat. avec produits finis

% L21–25, beneath J1 and J2
$J_1$ $=$ analyseur engendré par $J_0$ (analyseur \unsure{vierge|\deleted{libre}}) et les flèches données -- analyseur libre
\qquad $J_2$ $=$ sous-analyseur plein de $R$ engendré par $I\to R$, ens. sous-jacent $J_0$

% L19–25, lower right
\margin{$R$ apparait comme une catégorie de foncteurs de \deleted{\unsure{préfaisc}} $\widehat J_1$, où $J_1$ est l'analyseur libre sur $I$ engendré par les flèches induites (sur $J_0$) par $R$, $\widehat J_1$}

% L26–27, foot
\diagram{\begin{tikzcd}[ampersand replacement=\&,row sep=small,column sep=small]
 \widehat J_1\arrow[r,dash,"\mathrm{loc?}"]\arrow[d,"\mathrm{loc}"']\&\widehat J_2\arrow[d,"\mathrm{loc}"]\arrow[r,"\deleted{\unsure{\mathrm{loc}}}"]\&R\\
 R_1\arrow[r,dash,"\mathrm{loc}"]\&R_2\arrow[r,"\mathrm{loc}"]\&R_3
\end{tikzcd}}
% archive sheet number: 22
